\section{Conclusion}
\label{sec:idp-conclusion}

\EB{Não revisei esta seção ainda}

The experimental findings presented in this chapter underscore the critical importance of balancing workload distribution for efficient data-parallel processing in \ac{HPC} environments.
In essence, partitioning data into appropriately sized chunks is \emph{fundamental}: overly large chunks lead to \ac{OOM} failures, while excessive fragmentation forces undue scheduling overheads.
Achieving a balanced chunking strategy is therefore crucial for maximizing throughput and ensuring that computations proceed reliably, especially in domains such as seismic imaging and large-scale tensor-based analyses.

\vspace{6pt}
\noindent
\textbf{Predictive Load Balancing via Memory-Awareness.}
\begin{itemize}
    \item \emph{No \ac{OOM} Failures:} A central observation from our experiments is that \emph{memory-aware chunking}, underpinned by a predictive memory model, yielded \emph{zero} \ac{OOM} incidents. This reliability makes it possible to execute workloads safely even in tighter memory scenarios or with higher worker counts, where heuristic-based methods typically fail.
    \item \emph{Slight Speed Trade-Off:} While memory-aware chunking is typically slower than the fastest alternative (e.g., evenly-split or auto-chunking) on small-to-medium data, it remains fairly close to the best strategy in many scenarios. As dataset sizes increase, the ability to avoid \ac{OOM} errors becomes paramount, offsetting the extra runtime and positioning memaware as a robust fallback in virtually any setting.
    \item \emph{Predictability and Efficiency:} By predicting memory usage from input shapes and algorithm profiles, chunks can be sized so that the per-worker load stays below a configured memory limit. This approach dramatically reduces trial-and-error tuning, freeing up engineering effort while ensuring that computation remains safely within system resources.
\end{itemize}

\vspace{6pt}
\noindent
\textbf{Load Balancing Advantages.}
\begin{itemize}
    \item \textit{Resource Utilization:} An optimally balanced workload helps ensure workers operate near capacity without triggering memory overruns. In cluster contexts, such robust partitioning can prevent job crashes that might otherwise stall entire pipelines.
    \item \textit{Parallel Efficiency:} Even partitioning reduces idle times by preventing scenarios where one worker is overloaded while another finishes early. Memory-aware chunking further enforces a conservative boundary on chunk size, allowing high concurrency while staying safe from memory spikes.
\end{itemize}

\vspace{6pt}
\noindent
\textbf{Complexities and Future Directions.}
\begin{itemize}
    \item \textbf{Integration with Dask:}
    Although memory-aware chunking works within Dask’s framework, deeper integration could enable \emph{dynamic} adjustments to chunk sizes in response to real-time cluster metrics (e.g., CPU/GPU usage, memory availability).
    \item \textbf{Operator-Specific Extensions:}
    Our current memory model targets seismic and tensor-based computations. Adapting it to additional workloads (e.g., deep learning or PDE solvers) would involve gathering and training on a wider range of memory-consumption profiles.
    \item \textbf{Safety Factor Calibration:}
    Memory-aware chunking employs a user-defined safety factor to account for overhead (e.g., temporary arrays in intermediate steps). Research into an \emph{adaptive} approach could optimize resource usage further while remaining robust against \ac{OOM} risks.
    \item \textbf{Hybrid and Multi-Node Scenarios:}
    Scaling from a single \ac{HPC} node to large clusters—or across heterogeneous hardware—may require considering network bandwidth, \ac{I/O} patterns, and different memory topologies when predicting optimal chunk dimensions.
\end{itemize}

\vspace{6pt}
\noindent
\textbf{Implications for \ac{HPC} Workloads.}
By eliminating \ac{OOM} failures and maintaining predictable per-worker usage, memory-aware chunking enables workloads to scale more safely.
When data volumes become large, the slight overhead in runtime is often outweighed by the avoided failures and reduced idle time.
In practice, fewer failed jobs mean reduced operational costs and more reliable timetables for \ac{HPC} or cloud-based pipelines.

\vspace{6pt}
\noindent
\textbf{Concluding Remarks.}
In summary, memory-aware chunking offers a \emph{robust and data-driven} solution to the longstanding challenge of workload imbalance in parallel computing.
Not only does this method avoid \ac{OOM} failures in \emph{all} tested scenarios, but it also remains close to the performance of the best-performing chunking strategies in many cases.
Consequently, it stands as a powerful alternative for users seeking a well-rounded approach that balances stability, throughput, and scalability without the frustration of manual or trial-and-error fine-tuning.

Ultimately, while memory-aware chunking is not a cure-all for every computational challenge, it demonstrates that \emph{model-guided partitioning} can significantly improve reliability and throughput in data-parallel tasks.
Future work that integrates predictive modeling directly into task schedulers, expands its domain coverage, and refines safety-factor logic will only increase its versatility, opening new avenues for automated \ac{HPC} resource management and more intelligent data-parallel workflows.